\documentclass[11pt,a4paper]{article}
\usepackage[etiket=en]{gtu-el-yazisi}
\begin{document}
\notbilgi{Kontrol Sistemleri}{Positive feedback root locus, lead/lag compensators}{Control-Systems-lecture-12.pdf}

\sayfa{1}
\section{Positive Feedback architecture}

\elcizimi[0.7]{sekiller/Control-Systems-lecture-12-s1-1.png}[1]

\[
\frac{C(s)}{R(s)} = \frac{k\,G(s)}{1 - k\,G(s)H(s)}
\]
\[
1 - k\,G(s)H(s) = 0
\]
\[
G(s)H(s) = \overbrace{\frac{1}{k}}^{\text{positive}} + 0j
\]
\[
\angle G(s)H(s) = \tan^{-1}\!\left(\frac{0}{\varepsilon^{+}}\right) = \pm 360\,i \ \text{ or } \ \pm 180(2i)
\]
\formulkutu{|G(s)H(s)| = \frac{1}{k}}

positive feedback

\elcizimi[0.65]{sekiller/Control-Systems-lecture-12-s1-2.png}[1]

(2) Departure angles from \textit{complex} poles\\
arrival \quad '' \quad to \textit{complex} zero

\sayfa{2}
\noindent departure angle $= 360$ \okunamadi{} $+\sum$ phase of c.c pole w.r.t. other poles\\
\hspace*{5.2em}$-\sum$ phase of c.c pole w.r.t. other zeros

\medskip
\noindent (3) Asymptote angles
\[
\frac{\pm 360\,i}{n-m}
\]

\begin{ornek}
$G(s)=\dfrac{k\,(s+2)}{(s+3)(s^2+2s+2)}$\\
Positive feedback system
\end{ornek}

\begin{align*}
z_1 &= -2\\
p_1 &= -3\\
p_{2,3} &= -1 \pm j
\end{align*}

\elcizimi[0.55]{sekiller/Control-Systems-lecture-12-s2-1.png}[2]

\noindent branches: 3\\
$n-m = 3-1 = 2$\\
angle of asymptotes: $\dfrac{\pm 360\,i}{n-m} = \pm 180\,i$

\medskip
\noindent for $s=-1+j$
\[
\theta_1 = 0^\circ - 27^\circ - 90^\circ + 45^\circ = -72^\circ
\]
for $s=-1-j$
\[
\theta_2 = 0^\circ + 27^\circ + 90^\circ - 45^\circ = +72^\circ
\]

\medskip
\noindent To find characteristic equation
\[
1 + \frac{k\,(s+2)}{(s+3)(s^2+2s+2)} = 0
\]
\footnotetext{Claude notu: pozitif geri besleme için karakteristik denklem muhtemelen $1-kG(s)=0$ olmalı; $dk/ds=0$ kökleri bundan etkilenmez.}
\[
(s+3)(s^2+2s+2) + k\,(s+2) = 0
\]
\[
k = -\frac{(s+3)(s^2+2s+2)}{(s+2)}
\]
\begin{align*}
\frac{dk}{ds} = 0 &\;\Rightarrow\; 2s^3+11s^2+20s+10 = 0\\
&\;\Rightarrow\; 2(s+0{,}8)(s+2{,}35+j\,0{,}77)(s+2{,}35-j\,0{,}77) = 0
\end{align*}

\sayfa{3}
\noindent $\Rightarrow\ s=-0{,}8$ \ break in

\medskip
\noindent if its negative feedback

\elcizimi[0.5]{sekiller/Control-Systems-lecture-12-s3-1.png}[3]

\[
\text{centroid} = \frac{\sum \text{real part poles} - \sum \text{real part of zeroes}}{n-m} = -1{,}5
\]

\begin{ornek}
Negative feedback
\end{ornek}

\elcizimi[0.9]{sekiller/Control-Systems-lecture-12-s3-2.png}[3]

\noindent\textbf{Pozitive feedback}

\elcizimi[0.95]{sekiller/Control-Systems-lecture-12-s3-3.png}[3]

\noindent\textbf{Negative feedback}

\elcizimi[0.55]{sekiller/Control-Systems-lecture-12-s3-4.png}[3]

\begin{gather*}
\frac{\pm 180(2i+1)}{n-m}\\
\pm 45(2i+1)\\
\pm 45 \qquad \pm 135
\end{gather*}
centroid $=-2{,}5$

\sayfa{4}
\noindent Pozitif f\okunamadi{}

\elcizimi[0.65]{sekiller/Control-Systems-lecture-12-s4-1.png}[4]

\[
\frac{\pm 360\,i}{4} = \pm 90\,i
\]

\section{Compensator design}

\elcizimi[0.5]{sekiller/Control-Systems-lecture-12-s4-2.png}[4]

\noindent specific ``$k$'' which renders the system unstable

\noindent 4 seperate roots of ch.e.q.

\elcizimi[0.55]{sekiller/Control-Systems-lecture-12-s4-3.png}[4]

\[
1 + \frac{k\,(s+2)}{s^3\,(s+p)^3}
\]
\[
\theta = \cos^{-1}\zeta \qquad \text{damping ratio}
\]

\elcizimi[0.7]{sekiller/Control-Systems-lecture-12-s4-4.png}[4]

\noindent $\searrow\ \dfrac{1}{s+p_3} \rightarrow$ delayed integrator

\[
F(s) = \frac{s+a}{s+b} = \frac{R\,e^{i\phi}}{M\,e^{i\theta}} = \left(\frac{R}{M}\right) e^{i(\phi-\theta)}
\]

\noindent Phase of \okunamadi{s} points \quad $\beta = \phi_1 - \theta_1 - \theta_2$

\noindent the original OLTF \quad $\dfrac{s+z_1}{(s+p_1)(s+p_2)}$

\noindent if $\beta = 190^\circ$ \quad Negative \ $\dfrac{-10}{180}$ \ $\rightarrow$ to make it negative $p_3$ angle have to be bigger

\medskip
\noindent \okunamadi{PI} \quad $\dfrac{s+z_2}{s+p_3} \simeq \dfrac{1+z_2}{\lim_{p_3\to 0} s}$

\sayfa{5}
\noindent if $\beta = 150^\circ$ \quad $\dfrac{30^\circ}{180^\circ}$

\elcizimi[0.65]{sekiller/Control-Systems-lecture-12-s5-1.png}[5]

\[
\underbrace{(s+z_2)}_{\text{derivative operator}}\ \frac{(s+z_1)}{(s+p_1)(s+p_2)}
\]

\noindent Its unuseless itself so we have to add a pole

\elcizimi[0.7]{sekiller/Control-Systems-lecture-12-s5-2.png}[5]

\[
\frac{s+z_2}{s+p_3} \underset{\lim z_2\to 0}{=} \frac{s}{s+p_3} = 1 - \frac{p_3}{s+p_3}
\]
\noindent \textit{?? WTF}

\section{Adding poles}

\elcizimi[0.95]{sekiller/Control-Systems-lecture-12-s5-3.png}[5]

\noindent by adding poles the locus plot move to right, systems get more unstable. \hfill (lag compensator)

\sayfa{6}
\noindent \textcolor{kirmizikalem}{System get more unstable - slower! \ ($t_s\uparrow$)}

\noindent More accurate by adding integrator \\
\textcolor{kirmizikalem}{\underline{lag compensation,}} \ The output angle is less than input

\noindent \textcolor{kirmizikalem}{zero placement}

\elcizimi[0.56]{sekiller/Control-Systems-lecture-12-s6-1.png}[6]
\elcizimi[0.47]{sekiller/Control-Systems-lecture-12-s6-2.png}[6]
\elcizimi[0.58]{sekiller/Control-Systems-lecture-12-s6-3.png}[6]
\[ \frac{-1-3-(-2)}{2}=\frac{-4+2}{2}=-1 \]
\elcizimi[0.42]{sekiller/Control-Systems-lecture-12-s6-4.png}[6]

\noindent pulls the root locus plot to the left

\noindent $\rightarrow$ more stable

\noindent $\rightarrow$ faster response \quad ($t_s\downarrow$)

\noindent \hspace*{1.5em}more sensitive to noise

\noindent The output angle is more than the input angle

\noindent its leading output
\[ r(t)=\sin\omega t \quad\rightarrow\quad y(t)=\sin(\omega t+\phi) \]

\sayfa{7}
\elcizimi[0.42]{sekiller/Control-Systems-lecture-12-s7-1.png}[7]
\noindent \textcolor{kirmizikalem}{lag compensation} \ \okunamadi{$z>p>0$}

\noindent lead laging compensator \\
\textcolor{kirmizikalem}{Series/parallel compensation}

\elcizimi[0.85]{sekiller/Control-Systems-lecture-12-s7-2.png}[7]
\[ \frac{C(s)}{R(s)}=\frac{G_cG}{1+G_cGH} \]
\[ 1+G_cGH=0 \ \Rightarrow\ G_cGH=-1 \]
\[ \angle G_cGH = \text{gives the root locus} \]

\noindent parallel

\elcizimi[0.85]{sekiller/Control-Systems-lecture-12-s7-3.png}[7]
\[ \frac{C(s)}{R(s)}=\frac{G_1G_2}{1+G_2G_c+G_1G_2H} \qquad\longrightarrow\qquad 1+G_c\,\frac{G_2}{1+G_1G_2H}=0 \]
\[ 1+G_2G_c+G_1G_2H=0 \qquad\qquad \Rightarrow\ 1+G_cG_p=0 \]
\[ \times\ \frac{1}{1+G_1G_2H} \]
\[ =\frac{1+G_1G_2H+G_2G_c}{1+G_1G_2H} \]

\sayfa{8}
\section{Lead compensator: zero is closer to $j\omega$ axis}
\[ G_c(s)=k_c\,\frac{s+\frac{1}{T}}{s+\frac{1}{\alpha T}}\ \textcolor{turuncukalem}{\rightarrow\ \text{specs}} \qquad 0<\alpha<1 \]
\[ k_c\,\frac{s+z}{s+p} \]

\noindent if $\alpha>1\ \Rightarrow$ lag compensator
\[ G_c(s)=k_c\,\frac{s+\frac{1}{T}}{s+\frac{1}{\beta T}} \qquad \beta>1 \]
\[ \frac{1}{T}=10 \qquad \begin{cases} z=-10 \\ p=-\dfrac{10}{\alpha}=-20 \end{cases} \qquad \alpha=\frac{1}{2} \]

\elcizimi[0.8]{sekiller/Control-Systems-lecture-12-s8-1.png}[8]

\section{Lead lag compensator}
\elcizimi[0.75]{sekiller/Control-Systems-lecture-12-s8-2.png}[8]
\[ \left(k_c\,\frac{s+\frac{1}{T_1}}{s+\frac{1}{\alpha T_1}}\right)\left(\frac{s+\frac{1}{T_2}}{s+\frac{1}{\beta T_2}}\right) \qquad \text{\okunamadi{5 unknowns} to be design} \]
\[ 0<\alpha<1 \qquad\qquad \beta>1 \]

\noindent \textcolor{kirmizikalem}{Ex:}

\elcizimi[0.82]{sekiller/Control-Systems-lecture-12-s8-3.png}[8]

\sayfa{9}
\noindent Dominant pole $=-3+2j$ \quad \textcolor{turuncukalem}{$\rightarrow$ lead compensator}

\noindent \okunamadi{$e_{ss}$} $\leftarrow$ Steady state error $=10\%$ \ (step input) \quad \textcolor{turuncukalem}{$\rightarrow$ lag compensator}

\noindent $G_c=1$

\elcizimi[0.5]{sekiller/Control-Systems-lecture-12-s9-1.png}[9]
\[ G_c=\frac{s+\frac{1}{T}}{s+\frac{1}{\alpha T}} \]
\elcizimi[0.7]{sekiller/Control-Systems-lecture-12-s9-2.png}[9]
\[ \underline{\angle(-3+2j)} = -\theta_1-\theta_2+\phi = -\left(180-\tan^{-1}\!\left(\tfrac{2}{2}\right)\right)-90+\tan^{-1}\!\left(\tfrac{2}{1}\right) \]
\[ = -135-90+63{,}43^\circ = -161{,}56^\circ \]

\noindent phase deficiency $=-135-90+63{,}43^\circ=-161{,}56^\circ$

\noindent So compensator \quad $G_c(s)=\dfrac{s+4}{s+9}$
\[ E_{ss}=\frac{1}{1+k_p} \ \rightarrow\ \text{static position error constant} \]
\[ k_p=\lim_{s\to0}G(s) \]
\[ E_{ss}=\lim_{s\to0}\frac{sR(s)}{1+G(s)}=\lim_{s\to0}\frac{1}{1+G(s)}=\frac{1}{1+\lim_{s\to0}G(s)}=\frac{1}{1+k_p} \]

\sayfa{10}
\[ G(s)=k_c\left(\frac{s+4}{s+9}\right)\left(\frac{1}{(s+1)(s+3)}\right) \qquad \text{compensator system} \]
\[ E_{ss}=0{,}1=\frac{1}{1+k_p} \ \Rightarrow\ 1+k_p=10 \ \Rightarrow\ k_p=9 \]
\[ 9=k_p=\lim_{s\to0}G(s)=k_c\left(\frac{4}{9}\right)\left(\frac{1}{1\times3}\right) \]
\[ k_c\left(\frac{4}{27}\right)=9 \]
\[ \Rightarrow\ k_c=\frac{9\times27}{4} \]

\medskip\noindent\rule{0.45\linewidth}{0.4pt}\medskip

\[ G(s)=\left|\,k\,\frac{s+4}{s+9}\ \frac{1}{(s+1)(s+3)}\,\right|=1 \]

\noindent \hspace*{3em}$k$ \ give dominant pole

\end{document}
